% !TEX program = lualatex
% Transparencias de Fundamentos de las Matemáticas II
% Clase 10: Aritmética modular II

\documentclass[aspectratio=169,11pt]{beamer}

\usepackage{fontspec}
\usepackage[spanish,es-nodecimaldot]{babel}
\usepackage{amsmath,amssymb,mathtools}
\usepackage{booktabs}
\usepackage{csquotes}
\usepackage{xcolor}

\usetheme{Madrid}
\usecolortheme{default}
\setbeamertemplate{navigation symbols}{}
\setbeamertemplate{footline}[frame number]

\definecolor{FdMBlue}{RGB}{20,55,100}
\definecolor{FdMRed}{RGB}{130,30,30}
\definecolor{FdMGreen}{RGB}{25,100,60}

\setbeamercolor{structure}{fg=FdMBlue}
\setbeamercolor{title}{fg=white,bg=FdMBlue}
\setbeamercolor{frametitle}{fg=white,bg=FdMBlue}
\setbeamercolor{block title}{fg=white,bg=FdMBlue}
\setbeamercolor{block body}{bg=blue!3}
\setbeamercolor{alerted text}{fg=FdMRed}

\setmainfont{Latin Modern Roman}
\setsansfont{Latin Modern Sans}
\setmonofont{Latin Modern Mono}

\newcommand{\N}{\mathbb{N}}
\newcommand{\Z}{\mathbb{Z}}
\newcommand{\Q}{\mathbb{Q}}
\newcommand{\R}{\mathbb{R}}
\newcommand{\C}{\mathbb{C}}
\newcommand{\Pow}{\mathcal{P}}
\newcommand{\id}{\operatorname{id}}
\newcommand{\mcd}{\operatorname{mcd}}
\newcommand{\mcm}{\operatorname{mcm}}
\newcommand{\im}{\operatorname{Im}}

\title[FdM II -- Clase 10]{Aritmética modular II}
\subtitle{Operaciones e inversos modulares}
\author{Antonio Falcó Montesinos}
\institute{Universidad CEU Cardenal Herrera}
\date{Fundamentos de las Matemáticas II}

\begin{document}

\begin{frame}
  \titlepage
\end{frame}

\begin{frame}{Objetivos de la clase}
\begin{itemize}
  \item Definir suma y producto módulo \(n\).
  \item Comprender que las operaciones están bien definidas.
  \item Caracterizar inversos modulares.
\end{itemize}
\end{frame}

\begin{frame}{Mapa de la clase}
\tableofcontents
\end{frame}

\section{Operaciones}

\begin{frame}{Operaciones}
\begin{block}{Suma y producto}
\begin{itemize}
  \item \([a]_n+[b]_n=[a+b]_n\).
  \item \([a]_n[b]_n=[ab]_n\).
  \item Las operaciones se definen sobre clases, no sobre representantes aislados.
\end{itemize}
\end{block}
\end{frame}

\begin{frame}{Operaciones}
\begin{block}{Bien definidas}
\begin{itemize}
  \item Si se cambian representantes por congruentes, el resultado pertenece a la misma clase.
  \item Esta comprobación es necesaria para que la operación tenga sentido.
  \item La buena definición evita depender de una elección arbitraria.
\end{itemize}
\end{block}
\end{frame}

\begin{frame}{Operaciones}
\begin{block}{Ejemplo}
\begin{itemize}
  \item En \(\mathbb{Z}/5\mathbb{Z}\), \([3]+[4]=[7]=[2]\).
  \item \([3][4]=[12]=[2]\).
  \item Los cálculos se reducen usando restos.
\end{itemize}
\end{block}
\end{frame}

\section{Inversos}

\begin{frame}{Inversos}
\begin{block}{Definición}
\begin{itemize}
  \item Un inverso de \([a]_n\) es una clase \([b]_n\) con \([a]_n[b]_n=[1]_n\).
  \item Equivale a \(ab\equiv 1\pmod n\).
  \item No todas las clases no nulas tienen inverso.
\end{itemize}
\end{block}
\end{frame}

\begin{frame}{Inversos}
\begin{block}{Criterio}
\begin{itemize}
  \item \([a]_n\) tiene inverso si y solo si \(\mcd(a,n)=1\).
  \item La prueba usa la identidad de Bézout.
  \item El inverso puede encontrarse con Euclides extendido.
\end{itemize}
\end{block}
\end{frame}

\section{Profundización}

\begin{frame}{Profundización}
\begin{block}{Ejemplo}
\begin{itemize}
  \item En módulo \(7\), \([3]\) tiene inverso.
  \item Como \(3\cdot 5=15\equiv 1\pmod 7\), el inverso es \([5]\).
  \item En módulo \(8\), \([2]\) no tiene inverso porque \(\mcd(2,8)=2\).
\end{itemize}
\end{block}
\end{frame}

\begin{frame}{Profundización}
\begin{block}{Advertencia}
\begin{itemize}
  \item No se puede dividir por una clase sin comprobar que es invertible.
  \item Cancelar factores puede ser inválido si el factor no tiene inverso.
  \item La aritmética modular no se comporta siempre como un cuerpo.
\end{itemize}
\end{block}
\end{frame}

\begin{frame}{Profundización}
\begin{block}{Pregunta de control}
\begin{itemize}
  \item ¿Tiene inverso \([4]_{9}\)?
  \item ¿Tiene inverso \([6]_{15}\)?
  \item Justificar usando el \(\mcd\).
\end{itemize}
\end{block}
\end{frame}


\section{Detalle y práctica}

\begin{frame}{Detalle y práctica}
\begin{block}{Operar con clases}
\begin{itemize}
  \item Si \(a\equiv a'\pmod n\) y \(b\equiv b'\pmod n\), entonces \(a+b\equiv a'+b'\pmod n\).
  \item También se conserva el producto: \(ab\equiv a'b'\pmod n\).
  \item Estas propiedades permiten calcular dentro de \(\mathbb{Z}/n\mathbb{Z}\).
\end{itemize}
\end{block}
\end{frame}

\begin{frame}{Detalle y práctica}
\begin{block}{Actividad breve}
\begin{itemize}
  \item Calcular \(17+29\) módulo \(6\).
  \item Calcular \(17\cdot 29\) módulo \(6\).
  \item Comparar el resultado reduciendo antes y después de operar.
\end{itemize}
\end{block}
\end{frame}

\begin{frame}{Cierre}
\begin{itemize}
  \item Las operaciones módulo \(n\) se hacen sobre clases.
  \item La buena definición es imprescindible.
  \item La invertibilidad equivale a coprimalidad con el módulo.
\end{itemize}
\end{frame}

\begin{frame}{Trabajo recomendado}
\begin{itemize}
  \item Releer las secciones correspondientes del manual.
  \item Identificar definiciones, símbolos nuevos y enunciados principales.
  \item Preparar dudas y ejemplos para el seminario asociado.
\end{itemize}
\end{frame}

\end{document}
